% Gauss, Werke X/1 (1917), native leaves40–43, printed33–36.
\SourceLeaf{040}\NoSourceHead
\SourceGroup{IV}
\HistoricTitle{DRITTER BEWEIS DES FUNDAMENTALSATZES BEI DEN QUADRATISCHEN RESTEN IN EINER NEUEN EINKLEIDUNG.}
\HistoricCitation{[Aus Handbuch 21, Bg, S. 6—8.]}
\par\indent 1. Es seyn $m,M$ zwei positive ungerade [theilerfremde] Zahlen und $mM=\mu$. Wir bezeichnen
\begingroup\small
\[
\begin{array}{r@{\ }l@{\ }r@{\ }l@{\quad}l}
\text{durch }f&\text{den Inbegriff der}&\tfrac12(m-1)&\text{Zahlen}&1,2,3,\ldots,\tfrac12(m-1),\\
f'&&\tfrac12(m-1)&&\tfrac12(m+1),\tfrac12(m+3),\ldots,m-1,\\
F&&\tfrac12(M-1)&&1,2,3,\ldots,\tfrac12(M-1),\\
F'&&\tfrac12(M-1)&&\tfrac12(M+1),\tfrac12(M+3),\ldots,M-1,\\
\varphi&&\tfrac12(\mu-1)&&1,2,3,\ldots,\tfrac12(\mu-1),\\
\varphi'&&\tfrac12(\mu-1)&&\tfrac12(\mu+1),\tfrac12(\mu+3),\ldots,\mu-1.
\end{array}
\]
\endgroup
\noindent Die Zahlen $f',F',\varphi'$ sind folglich identisch mit $m-f,\ M-F,\ \mu-\varphi$.

2. Wir bezeichnen ferner durch $A,B,C,D$ die Inbegriffe der kleinsten positiven Reste, welche nach dem Modulus $\mu$ resp. alle Zahlen
\[
fM+Fm,\ fM+F'm,\ f'M+Fm,\ f'M+F'm
\]
\noindent geben. Die Anzahl der Zahlen in einem jedem wird offenbar $\tfrac14(m-1)(M-1)$, also die Anzahl aller $(m-1)(M-1)$. Ferner sieht man leicht, dass alle unter sich ungleich und keiner durch $m$ noch durch $M$ theilbar seyn werden; woraus man den Schluss zieht, dass alle zusammen mit dem Inbegriffe aller weder durch $m$ noch durch $M$ theilbaren Zahlen zwischen den Grenzen $0$ und $\mu$, d. i. in den beiden Reihen $\varphi$ und $\varphi'$ identisch sind.
\vfill\noindent\small X$_1$.\hfill5\par\normalsize
\clearpage

\SourceLeaf{041}\ArithmeticHead{34}{L}{ARITHMETIK.\quad NACHLASS.}
\par\indent Noch bemerken wir, dass alle Zahlen $f'M+F'm$ identisch sind mit
\[
(m-f)M+(M-F)m
\]
\noindent d. i. mit
\[
2\mu-(fM+Fm);
\]
\noindent die Zahlen $D$ werden also die Ergänzungen der Zahlen $A$ zu $\mu$, und aus ähnlichem Grunde die Zahlen $C$ die Ergänzungen der Zahlen $B$ zu $\mu$ seyn.

3. Wir scheiden die Zahlen $A,B,C,D$ resp. in zwei Classen $\alpha$ und $\alpha'$, $\beta$ und $\beta'$, $\gamma$ und $\gamma'$, $\delta$ und $\delta'$, so dass in $\alpha,\beta,\gamma,\delta$ diejenigen fallen, welche kleiner als $\tfrac12\mu$, in $\alpha',\beta',\gamma',\delta'$ hingegen diejenigen, welche grösser als $\tfrac12\mu$ sind. Zusammengenommen werden also $\alpha,\beta,\gamma,\delta$ alle weder durch $m$ noch durch $M$ theilbaren Zahlen in $\varphi$, so wie $\alpha',\beta',\gamma',\delta'$ zusammen die ähnlichen Zahlen in $\varphi'$ darstellen. Da nun offenbar $Fm$ alle durch $m$ theilbaren Zahlen in $\varphi$ und $fM$ alle durch $M$ theilbaren in $\varphi$ gibt, und alle diese Zahlen ungleich seyn müssen (keine zugleich durch $m$ und $M$ theilbar seyn kann); so besteht $\varphi$ offenbar aus $\alpha,\beta,\gamma,\delta,Fm$ und $fM$; und auf ähnliche Weise $\varphi'$ aus $\alpha',\beta',\gamma',\delta'$, $F'm$ und $f'M$.

4. Die Anzahl der Zahlen, welche von $A,B,C,D$ in die Classen $\alpha,\beta,\gamma,\delta$ kommen, bleibt zwar noch unbestimmt; allein so viel ist klar, da $\delta$ und $\delta'$ zusammen die Ergänzungen von $\alpha$ und $\alpha'$ zusammen zu $\mu$ sind, dass nothwendig $\delta$ die Ergänzungen von $\alpha'$ und $\delta'$ die Ergänzungen von $\alpha$ seyn werden. Eben so müssen $\gamma,\gamma'$ resp. die Ergänzungen von $\beta'$ und $\beta$ zu $\mu$ seyn. Bedeutet also $p$ die Anzahl der Zahlen in $\alpha$, $q$ die Anzahl in $\beta$, so sind
\[
\begin{array}{ll}
\text{in }\delta'\text{ gleichfalls}&p\\
\text{in }\gamma'\text{ gleichfalls}&q\\
\text{in }\alpha'\text{ oder in }\delta\text{ hingegen}&\tfrac14(m-1)(M-1)-p\\
\text{in }\beta'\text{ oder in }\gamma\quad\text{aber}&\tfrac14(m-1)(M-1)-q.
\end{array}
\]
\par\indent 5. Die kleinsten Reste der $\tfrac12(m-1)$ Zahlen $fM$ nach dem Modulus $m$ liegen zum Theil in $f$, zum Theil in $f'$; die Anzahl der letztern sey $n$, also die Anzahl der erstern $\tfrac12(m-1)-n$. Eben so liegen die kleinsten Reste der $\tfrac12(M-1)$ Zahlen $Fm$ nach dem Modulus $M$ zum Theil in $F$, zum Theil in $F'$; die Anzahl der letztern $=N$ gesetzt, wird die der erstern $=\tfrac12(M-1)-N$ seyn.
\clearpage

\SourceLeaf{042}\ArithmeticHead{35}{R}{DRITTER BEWEIS DES FUNDAMENTALSATZES.}
\par\indent 6. Zählen wir jetzt alle Zahlen in $\varphi'$ zusammen, die nach dem Modulus $m$ einer der Zahlen $fM$ congruent sind. In so fern $\varphi'$ aus $\alpha',\beta',\gamma',\delta',f'M$, $F'm$ zusammengesetzt ist, und alle Zahlen $\alpha',\beta'$ offenbar die gedachte Eigenschaft haben und keine der übrigen sie hat, wird offenbar die verlangte Anzahl $=\tfrac12(m-1)(M-1)-p-q$ seyn. Offenbar besteht aber $\varphi'$ auch andererseits aus den Zahlen
\[
\begin{array}{lll}
\tfrac12(\mu+1),&\tfrac12(\mu+3)&,\ldots,\tfrac12(M+1)m-1\\
\tfrac12(M+1)m,&\tfrac12(M+1)m+1,&\ldots,\tfrac12(M+3)m-1\\
\tfrac12(M+3)m,&\tfrac12(M+3)m+1,&\ldots,\tfrac12(M+5)m-1\\
\tfrac12(M+5)m,&\tfrac12(M+5)m+1,&\ldots,\tfrac12(M+7)m-1\\
\multicolumn{3}{c}{\text{etc. bis}}\\
(M-1)m,&(M-1)m+1,&\ldots,\mu-1
\end{array}
\]
\noindent oder
\begingroup\small
\[
\begin{array}{l@{\quad}c}
&\begin{gathered}\text{unter welchen sich Zahlen}\\\text{von gedachter Eigenschaft}\\\text{befinden}\end{gathered}\\
\tfrac12(M-1)m+f'&n\\
\tfrac12(M+1)m,\ \tfrac12(M+1)m+f,\ \tfrac12(M+1)m+f'&\tfrac12(m-1)\\
\tfrac12(M+3)m,\ \tfrac12(M+3)m+f,\ \tfrac12(M+3)m+f'&\tfrac12(m-1)\\
\tfrac12(M+5)m,\ \tfrac12(M+5)m+f,\ \tfrac12(M+5)m+f'&\tfrac12(m-1)\\
\multicolumn{1}{c}{\text{etc.}}&\\
(M-1)m,\quad(M-1)m+f,\quad(M-1)m+f',&\tfrac12(m-1).
\end{array}
\]
\endgroup
\par\noindent\hspace*{\parindent}Auf diese Weise beträgt also die Anzahl aller $n+\tfrac14(M-1)(m-1)$ und man hat also die Gleichung
\[
n=\tfrac14(m-1)(M-1)-p-q.
\]
\par\indent 7. Scheiden wir auf ähnliche Weise aus $\varphi'$ alle Zahlen aus, welche nach dem Modulus $M$ einer der Zahlen $Fm$ congruent sind, so sind dies die Zahlen $\alpha'$ und $\gamma'$, also ihre Anzahl $\tfrac14(m-1)(M-1)-p+q$. Dieselbe Anzahl findet sich aber durch eine zweite Zerlegung von $\varphi'$ in
\[
\tfrac12M(m-1)+F',\ f'M,\ f'M+F,\ f'M+F'
\]
\noindent $=N+\tfrac14(M-1)(m-1)$. Wir haben also\par
\[
N=-p+q.
\]
\par\vfill\noindent\hfill5*\par
\clearpage

\SourceLeaf{043}\ArithmeticHead{36}{L}{ARITHMETIK.\quad NACHLASS.}
\par\indent 8. Es wird demnach
\[
N-n=2q-\tfrac14(m-1)(M-1);
\]
\noindent also sind $N$ und $n$ beide gerade oder beide ungerade, so oft $\tfrac14(m-1)(M-1)$ gerade, d. i. so oft unter den Zahlen $m,M$ wenigstens eine von der Form $4k+1$ ist; hingegen werden $N,n$ nothwendig die eine gerade, die andere ungerade seyn müssen, wenn $\tfrac14(m-1)(M-1)$ ungerade ist, d. i. wenn beide Zahlen $m,M$ zugleich von der Form $4k+3$ sind.

\textit{Hieraus folgt das Fundamentaltheorem von selbst.}

Bei dieser Einkleidung des Beweises ist der wahre Nerf desselben mehr in die Augen fallend als bei derjenigen, in welcher er in den Göttingischen Commentationen Bd. XVI\begingroup\renewcommand{\thefootnote}{[*)}\footnote{1808, Werke II, S. 1 ff.]}\textsuperscript{]}\endgroup\ erscheint.
\begin{flushright}Novbr. 12. [1813]\end{flushright}
\begin{center}\textbf{BEMERKUNG.}\end{center}
\begingroup\small
\par\indent Die vorstehende »neue Einkleidung« des dritten Beweises (Werke II, S. 1) ist (vergl. den Artikel 20 des \textsc{Bachmann}schen Aufsatzes »Über \textsc{Gauss}' zahlentheoretische Arbeiten«) nur in der Entwicklung von dem fünften Beweise (1818, Werke II, S. 51 ff.) verschieden.
\begin{flushright}\textsc{Schlesinger.}\end{flushright}
\ClosingRule\clearpage\endgroup
